\subsubsection{保証・セキュリティ・安全性}

\term{保証}{Assurance} は、ソフトウェアが意図どおりに動くという信頼を得ることに関わる。\term{セキュリティ}{Security} は、情報や資源を不正なアクセス、変更、破壊、利用不能化から守ることである。\term{安全性}{Safety} は、人命、財産、環境への危害を避けることである。

セキュリティ設計では、機密性、完全性、可用性、認証、認可、監査、攻撃時の被害限定、復旧を考える。安全性設計では、危険状態に至る条件、停止方法、警告、フェイルセーフ、冗長化を考える。どちらも後付けが難しいため、設計段階から扱う必要がある。
